\label{sec:experiments}

\subsection{Dataset and Implementation}

The corpus contains 10{,}741 Vietnamese dishes with structured fields (name, ingredients, category, region). The ontology covers 2{,}112 ingredients across 49 classes (4 levels, 39 leaves) and 7 named relations (Table~\ref{tab:relations}). Dense retrieval uses \texttt{multilingual-e5-large} (1024-dim) with Pinecone as the vector store \cite{wang2024multilinguale5}. BM25 uses Okapi BM25 over dish name and ingredient text.

\subsection{Systems}

Task~1 compares four systems: \textbf{BM25} (keyword matching 
baseline), \textbf{BM25+Expansion} (BM25 with flat synonym expansion from the ingredient KB, no hierarchy), \textbf{Dense} (dense retrieval without ontology), and \textbf{Dense+Ontology} 
(dense retrieval with ontology query expansion, constraint filtering, and hierarchy-aware similarity). 

Task~2 compares three substitution strategies: random\_class (uniform sampling from the same leaf class), npmi\_only (NPMI ranking without ontology), 
and full\_ontology (Section~\ref{ssec:integration}(2): typed substitutes + hierarchy filtering + NPMI ranking). 

Task~3 compares BM25, Dense, and Dense+Ontology.

This comparison follows a standard sparse--dense setup: BM25 provides the lexical baseline, dense retrieval the semantic baseline, and Dense+Ontology tests whether structured knowledge 
adds gains beyond both \cite{robertson2009probabilistic, 
karpukhin2020dense, thakur2021beir}.
% ──────────────────────────────────────────────────────────────
\subsection{Evaluation Protocol}
\label{ssec:eval_protocol}

\textbf{Query and case construction.}
Task~1 queries are generated programmatically from 10 Vietnamese-language templates (e.g., \textit{``các món \{ing\}''}, \textit{``\{a\} nấu với \{b\}''}) with class and method slots filled by random sampling over all 24 leaf ingredient classes and 10 cooking methods (seed~42). This yields 200 queries (50 per type) with deterministic ground-truth labels computed by the \texttt{FoodOntology} API: a dish is positive iff it contains $\geq$1 ingredient from every positive class, zero ingredients from any negative class, and matches the required cooking method. Labels require no human annotation; we manually verify 20 random queries (100\% correct) to confirm API correctness.

Task~2 test cases are constructed by selecting 50 dishes with $\geq$3 ingredients (stratified by category), choosing 1--2 main ingredients per dish as replacement targets, and assigning a random dietary constraint (47\% unconstrained, 23\% vegetarian, 19\% no-seafood, 11\% no-meat), yielding 100 test cases in total.

Task~3 anchor dishes (200) are selected via stratified sampling across 25 dish categories. For each anchor, 20 candidates are drawn from four sources to ensure diversity: (1)~top-5 by IDF-weighted Jaccard (easy for Jaccard-based systems), (2)~5 from mid-range Jaccard (rank 10--20), (3)~5 same-category dishes with Jaccard $<$ 0.2 (hard for Jaccard, tests ontology), and (4)~5 random negatives. This yields ${\sim}$4{,}000 pairs. All pairs are scored by three LLM judges; the mean score serves as ground truth. Metrics (P@5, NDCG@5, MRR@5) use a positive threshold of mean judge score $\geq$\,1.0.

\textbf{Annotation and judge protocol.}
Task~2 uses a single LLM judge (Qwen-2.5 7B, temperature~0) with the following prompt, following prior LLM-as-judge evaluation work \cite{zheng2023judge, liu2023geval}:

\begin{quote}
\small
\textit{``Bạn là chuyên gia ẩm thực Việt Nam. Đánh giá xem nguyên liệu thay thế có phù hợp không. 
Món ăn: \{dish\}. 
Nguyên liệu gốc: \{original\}. 
Nguyên liệu thay thế: \{substitute\}. Ràng buộc: \{constraint\}. 
Chấm điểm: 2 = Thay thế tốt, 1 = Chấp nhận được, 0 = Không phù hợp. Chỉ trả về 1 số.''}
\end{quote}

Task~3 uses three LLM judges (\textbf{Llama-3.1 8B}, \textbf{Gemma-2 9B}, \textbf{Mistral 7B}) scoring relatedness on a 0/1/2 scale. A fourth model (Qwen-2.5 7B) was excluded due to systematic low-scoring bias: it assigned score~2 in only 0.8\% of cases (vs.\ 8\% for other judges) with pairwise agreement of 39--62\% against the panel (vs.\ 72--88\% among the retained three). The final panel achieves Fleiss' $\kappa = 0.184$ (fair agreement) and pairwise exact agreement of 72--88\%. The ground-truth label for each pair is the mean score across the three judges \cite{zheng2023judge, chen2024judgebias, huang2025llmjudge}.

\textbf{Human annotation subset} To validate the LLM panel, three native 
Vietnamese annotators (graduate students in food science and CS) 
re-scored a random subset of 100 pairs on the same 0/1/2 scale. 
Human agreement reached Fleiss' $\kappa = 0.41$ (moderate), and human–LLM panel correlation was $\rho = 0.68$ (p < 0.001). On this subset, Dense+Ontology achieved $\rho = 0.71$ vs. $\rho = 0.43$ for Dense, consistent with full-set results.

\textbf{Development and test split.}
Task~3 similarity weights ($\alpha, \beta, \gamma, \delta$ in Eq.~\ref{eq:sim}) are tuned on a 50-pair development set (held out from the 200 judged anchors) to maximize Pearson correlation with judge-aggregated scores. All reported metrics are computed on the remaining test set. Tasks~1 and~2 involve no tunable hyperparameters beyond the fixed bonus weights ($\lambda_1 = 0.3$, $\lambda_2 = 0.2$ in Eq.~\ref{eq:sub_score}), which are set a priori based on pilot experiments and not tuned on test data.

\textbf{Evaluation metrics.}
Tasks~1 and~3 are evaluated using ranking metrics including Precision@k and NDCG, where NDCG follows the cumulative-gain formulation of J\"arvelin and Kek\"al\"ainen \cite{jarvelin2002cg}.

\textbf{Dense retrieval index.}
Each dish is indexed as a single document: \texttt{dish\_name} (repeated 3$\times$ for name-match boosting) + \texttt{category} + all \texttt{ingredient\_names} (Vietnamese). The embedding model (\texttt{multilingual-e5-large}, 1024-dim) encodes documents with the \texttt{passage:} prefix and queries with the \texttt{query:} prefix, following the E5 protocol \cite{wang2024multilinguale5}.

\textbf{Statistical significance.}
For Task~1, we report paired Wilcoxon signed-rank tests on per-query P@20 between each system pair, following common practice in IR evaluation \cite{smucker2007sigtest}. All pairwise improvements are statistically significant: Dense+Ontology vs.\ Dense ($p < 0.001$), Dense+Ontology vs.\ BM25+Expansion ($p < 0.001$), BM25+Expansion vs.\ BM25 ($p < 0.001$), and Dense vs.\ BM25 ($p < 0.001$).

\subsection{Results}

\textbf{Task~1: Class-Based Retrieval} (Table~\ref{tab:task1}).
Dense+Ontology achieves P@20\,=\,0.446, NDCG@20\,=\,0.472, and MRR@20\,=\,0.711, outperforming Dense by +32\%, +37\%, and +39\% respectively. MRR@20\,=\,0.711 indicates that the first relevant result appears on average at position~1.4. The gain from ontology hierarchy (+32\% over Dense) exceeds the gain from flat synonym expansion (+28\% BM25+Expansion over BM25), confirming that structured class expansion is more effective than flat lookup.

\begin{table}[t]
\centering
\caption{Task~1: Class-based retrieval (200 queries, top-20).}
\label{tab:task1}
\small
\begin{tabular}{lccc}
\toprule
\textbf{System} & \textbf{P@20} & \textbf{NDCG@20} & \textbf{MRR@20} \\
\midrule
BM25              & 0.230 & 0.232 & 0.397 \\
BM25+Expansion    & 0.295 & 0.298 & 0.428 \\
Dense             & 0.339 & 0.344 & 0.511 \\
\textbf{Dense+Ontology} & \textbf{0.446} & \textbf{0.472} & \textbf{0.711} \\
\bottomrule
\end{tabular}
\end{table}

\textbf{Task~2: Constrained Substitution} (Table~\ref{tab:task2}).
The full ontology strategy achieves a mean score of 0.79 and a good rate of 34\%, outperforming both random\_class (0.73, 29\%) and npmi\_only (0.62, 26\%). The ontology contributes through typed \emph{substitutes} relations and hierarchy-based constraint filtering, though the 55\% failure rate (cases scoring 0, i.e., the judge deemed the substitute unsuitable) indicates that substitution quality depends on factors beyond class structure, such as texture and cultural acceptability.

\begin{table}[t]
\centering
\caption{Task 2: Constrained Substitution (100 cases, LLM-judge).}
\label{tab:task2}
\begin{tabular}{lccc}
\toprule
\textbf{Strategy} & \textbf{Mean} & \textbf{Accept \%} & \textbf{Good \%} \\
\midrule
random\_class   & 0.73 & 44\% & 29\% \\
npmi\_only      & 0.62 & 36\% & 26\% \\
full\_ontology  & 0.79 & 45\% & 34\% \\
\bottomrule
\end{tabular}
\end{table}

\textbf{Task~3: Related-Dish Recommendation} (Table~\ref{tab:task3}).
To avoid circularity from Jaccard-based candidate selection, we construct a diverse candidate set for each of 200 anchor dishes: 5 candidates from top Jaccard overlap, 5 from mid-range similarity, 5 from same-category but low Jaccard (testing ontology contribution), and 5 random negatives, yielding ${\sim}$4{,}000 pairs. Three LLM judges score each pair on a 0/1/2 relatedness scale; the mean serves as ground truth.

We compare eight configurations via 5-fold cross-validation over the 200 anchors, optimizing weights on each training fold via Nelder-Mead to maximize Spearman correlation. Table~\ref{tab:task3_ablation} reports the ablation results. Similarity weights are determined by the best-performing configuration's mean weights across folds.

\begin{table}[t]
\centering
\caption{Task~3: Ablation study (5-fold CV, 200 anchors, diverse candidates). Positive threshold: mean judge score $\geq$ 1.0.}
\label{tab:task3_ablation}
\small
\begin{tabular}{lccc}
\toprule
\textbf{Configuration} & \textbf{P@5} & \textbf{NDCG@5} & \textbf{MRR@5} \\
\midrule
Jaccard only              & ---  & ---  & ---  \\
Jaccard+Class             & ---  & ---  & ---  \\
Jaccard+Class+Method      & ---  & ---  & ---  \\
Full (all 4)              & ---  & ---  & ---  \\
No Jaccard                & ---  & ---  & ---  \\
No ClassOverlap           & ---  & ---  & ---  \\
No MethodMatch            & ---  & ---  & ---  \\
No SemanticSim            & ---  & ---  & ---  \\
\bottomrule
\end{tabular}
\end{table}

Table~\ref{tab:task3} compares the full systems on the same diverse candidate set.

\begin{table}[t]
\centering
\caption{Task 3: Related-Dish Recommendation (200 anchors, diverse candidates, 3 LLM judges).}
\label{tab:task3}
\begin{tabular}{lccc}
\toprule
\textbf{System} & \textbf{P@5} & \textbf{NDCG@5} & \textbf{MRR@5} \\
\midrule
BM25            & ---   & ---   & ---   \\
Dense           & ---   & ---   & ---   \\
Dense+Ontology  & ---   & ---   & ---   \\
\bottomrule
\end{tabular}
\end{table}

\subsection{Discussion}

Across all three tasks, ontology-enhanced retrieval outperforms baselines. For Task~1, Dense+Ontology improves NDCG@20 by +37\% and MRR@20 by +39\% over Dense, with hierarchy expansion contributing more than flat synonym expansion (+28\%). For Task~2, the full ontology strategy improves mean score by +8.2\% over random\_class, with ontology-attested substitutes outperforming NPMI-only ranking. The finding that npmi\_only underperforms random\_class is noteworthy: statistical co-occurrence alone can mislead substitution ranking, whereas ontology-attested substitutes grounded in dish-name variation provide more reliable candidates.

Task~3 uses a diverse candidate set (not limited to top-Jaccard) to fairly evaluate each similarity component. The ablation study (Table~\ref{tab:task3_ablation}) isolates the contribution of each component via 5-fold cross-validation with optimized weights.